\documentclass[11pt]{article}

% Packages
\usepackage[margin=1in]{geometry}
\usepackage{listings}
\usepackage{xcolor}
\usepackage{hyperref}

% Python code styling
\lstset{
    language=Python,
    basicstyle=\ttfamily\small,
    keywordstyle=\color{blue},
    commentstyle=\color{gray},
    stringstyle=\color{red},
    numbers=left,
    numberstyle=\tiny\color{gray},
    stepnumber=1,
    numbersep=5pt,
    backgroundcolor=\color{white},
    showspaces=false,
    showstringspaces=false,
    showtabs=false,
    frame=single,
    tabsize=4,
    captionpos=b,
    breaklines=true,
    breakatwhitespace=false,
}

\title{Sakthy re-learns Python as Senior}

\begin{document}

\maketitle
\tableofcontents
\newpage

\section{Motivation}
Python is a high-level, interpreted programming language.

\section{Basic Syntax}

\subsection{Hello World}
\begin{lstlisting}
print("Hello, World!")
\end{lstlisting}

\subsection{Variables}
\begin{lstlisting}
# Variables in Python
name = "Alice"
age = 25
height = 5.6
is_student = True
\end{lstlisting}

\section{Data Structures}

\subsection{Lists}
\begin{lstlisting}
# Creating a list
fruits = ["apple", "banana", "cherry"]
print(fruits[0])  # Output: apple
\end{lstlisting}

\subsection{Dictionaries}
\begin{lstlisting}
# Creating a dictionary
person = {
    "name": "Bob",
    "age": 30,
    "city": "New York"
}
print(person["name"])  # Output: Bob
\end{lstlisting}

\section{Control Flow}

\subsection{If Statements}
\begin{lstlisting}
x = 10
if x > 5:
    print("x is greater than 5")
elif x == 5:
    print("x is 5")
else:
    print("x is less than 5")
\end{lstlisting}

\subsection{Loops}
\begin{lstlisting}
# For loop
for i in range(5):
    print(i)

# While loop
count = 0
while count < 5:
    print(count)
    count += 1
\end{lstlisting}

\section{Functions}
\begin{lstlisting}
def greet(name):
    """Function to greet a person"""
    return f"Hello, {name}!"

# Call the function
message = greet("Alice")
print(message)
\end{lstlisting}

\end{document}
